%----------------------------------------------------------------------------------------
% METADATA BLOCK (YAML)  — see inbox/README.md
%----------------------------------------------------------------------------------------
% uid: ""
% school: "lycee"
% level: "2eme"
% section: ["Sciences", "Informatique"]
% chapter: "Notion de Polynômes"
% title: "Racine évidente et factorisation"
% difficulty: "Facile"
% tags: ["polynomes", "racine-evidente", "factorisation", "identification"]
% provenance: "own"
% batch: "2026-11-19-2sc-polynomes-barycentre"
%----------------------------------------------------------------------------------------

\newcommand{\R}{\mathbb{R}}
\newcommand{\vect}[1]{\overrightarrow{#1}}

% ---STATEMENT---
Pour chaque polynôme, chercher une racine parmi $-2$, $-1$, $1$ et $2$, factoriser le polynôme au maximum, puis résoudre dans $\R$ l'équation $P(x)=0$.
\begin{enumerate}
  \item $P_1(x)=x^3-6x^2+11x-6$ ;
  \item $P_2(x)=2x^3+x^2-5x+2$ ;
  \item $P_3(x)=3x^3-4x^2-5x+2$ ;
  \item $P_4(x)=x^3-2x^2+2x-1$.
\end{enumerate}

% ---SOLUTION---
\begin{prop}[Factorisation]
Si \ib{$P(\alpha)=0$}, alors $P(x)=(x-\alpha)\,Q(x)$ où $Q$ est un polynôme de degré $\deg P-1$, que l'on trouve par identification.
\end{prop}
\begin{prop}[Trinôme]
Pour $ax^2+bx+c$ : \ib{$\Delta=b^2-4ac$}. Si $\Delta>0$ : deux racines $\dfrac{-b\pm\sqrt\Delta}{2a}$ ; si $\Delta=0$ : une racine double $-\dfrac{b}{2a}$ ; si $\Delta<0$ : pas de racine.
\end{prop}
\begin{enumerate}[itemsep=14pt]
  \item $P_1(1)=1-6+11-6=0$ : $1$ est racine. On cherche $P_1(x)=(x-1)(ax^2+bx+c)=ax^3+(b-a)x^2+(c-b)x-c$ :
  \[
  \left\{
  \begin{aligned}
    a&=1\\
    b-a&=-6\\
    c-b&=11\\
    -c&=-6
  \end{aligned}
  \right.
  \qquad\text{donc}\qquad a=1,\ b=-5,\ c=6\quad(\text{contrôle : }6-(-5)=11\ \checkmark)
  \]
  \[x^2-5x+6:\quad\Delta=25-24=1,\qquad x=\dfrac{5-1}{2}=2\quad\text{ou}\quad x=\dfrac{5+1}{2}=3\]
  \[P_1(x)=(x-1)(x-2)(x-3)\qquad\mathcal S=\{1\,;2\,;3\}\]

  \item $P_2(1)=2+1-5+2=0$ : $1$ est racine. $P_2(x)=(x-1)(ax^2+bx+c)=ax^3+(b-a)x^2+(c-b)x-c$ :
  \[
  \left\{
  \begin{aligned}
    a&=2\\
    b-a&=1\\
    c-b&=-5\\
    -c&=2
  \end{aligned}
  \right.
  \qquad\text{donc}\qquad a=2,\ b=3,\ c=-2\quad(\text{contrôle : }-2-3=-5\ \checkmark)
  \]
  \[2x^2+3x-2:\quad\Delta=9+16=25,\qquad x=\dfrac{-3-5}{4}=-2\quad\text{ou}\quad x=\dfrac{-3+5}{4}=\dfrac12\]
  \[P_2(x)=(x-1)(x+2)(2x-1)\qquad\mathcal S=\left\{-2\,;\dfrac12\,;1\right\}\]

  \item $P_3(1)=-4\neq0$, mais $P_3(-1)=-3-4+5+2=0$ : $-1$ est racine. $P_3(x)=(x+1)(ax^2+bx+c)=ax^3+(a+b)x^2+(b+c)x+c$ :
  \[
  \left\{
  \begin{aligned}
    a&=3\\
    a+b&=-4\\
    b+c&=-5\\
    c&=2
  \end{aligned}
  \right.
  \qquad\text{donc}\qquad a=3,\ b=-7,\ c=2\quad(\text{contrôle : }-7+2=-5\ \checkmark)
  \]
  \[3x^2-7x+2:\quad\Delta=49-24=25,\qquad x=\dfrac{7-5}{6}=\dfrac13\quad\text{ou}\quad x=\dfrac{7+5}{6}=2\]
  \[P_3(x)=(x+1)(3x-1)(x-2)\qquad\mathcal S=\left\{-1\,;\dfrac13\,;2\right\}\]

  \item $P_4(1)=1-2+2-1=0$ : $1$ est racine. $P_4(x)=(x-1)(ax^2+bx+c)=ax^3+(b-a)x^2+(c-b)x-c$ :
  \[
  \left\{
  \begin{aligned}
    a&=1\\
    b-a&=-2\\
    c-b&=2\\
    -c&=-1
  \end{aligned}
  \right.
  \qquad\text{donc}\qquad a=1,\ b=-1,\ c=1\quad(\text{contrôle : }1-(-1)=2\ \checkmark)
  \]
  \[x^2-x+1:\quad\Delta=1-4=-3<0\quad\text{: pas de racine}\]
  \[P_4(x)=(x-1)(x^2-x+1)\qquad\mathcal S=\{1\}\]
  \textbf{Piège :} un polynôme de degré $3$ n'a pas toujours trois racines ; ici le trinôme ne se factorise pas.
\end{enumerate}
